\extrachap{Conventions}

Unless otherwise stated, $X$ is a connected compact complex manifold of complex dimension $n\geq1$, without boundary. A Hermitian form is a smooth real strictly positive $(1,1)$-form. We reserve $\omega$ for such a form and allow $\theta$ to be an arbitrary smooth real $(1,1)$-form. Closedness of $\theta$ or $\omega$ will always be stated.

We use the normalizations of \cite{SGPT}:
\[
 \dc=\frac{\mathrm i}{4\pi}(\bar\partial-\partial),
 \qquad \ddc=\frac{\mathrm i}{2\pi}\partial\bar\partial.
\]
On a disc in $\mathbb C$ with coordinate $z$, $\ddc\log|z|^2$ is the unit Dirac measure at the origin in complex dimension one. For a quasi-plurisubharmonic weight $u$, multiplier ideals use $\exp(-u)$, not $\exp(-2u)$. If $\lfloor\lambda\rfloor$ denotes the greatest integer at most a real number $\lambda$, for $\lambda\geq0$ this choice makes the ideal of $\lambda\log|z|^2$ equal to $(z^{\lfloor\lambda\rfloor})$.

The differential operator and integration differentials are set upright: $\mathrm{d}u$, $\mathrm{d}z$ and $\mathrm{d}\mu$. A letter $d$ denoting a dimension, degree or distance remains italic.

For quasi-plurisubharmonic functions $u,v$ on the same space, the notation $u\preceq v$ means $u\leq v+C$ for some real constant $C$: $u$ is more singular than $v$. We write $u\sim v$ when both inequalities hold. Inequalities between extended-real functions are used instead of taking their difference on a common pole set.

A quasi-plurisubharmonic function is locally a sum of a plurisubharmonic function and a smooth real function, and is not identically $-\infty$. We use its canonical upper semicontinuous representative. Positivity of a current is understood in the usual weak sense. For a smooth potential $u$, we write
\[
 \theta_u=\theta+\ddc u.
\]
Products involving nonsmooth potentials require a separate definition; smooth differential-form calculations do not by themselves justify such products.

All mass formulas below use $\theta_u^n$, without a factor $1/n!$. A numerical conversion is required when a reference uses a different normalization of $\ddc$ or of volume.

Each section and subsection declares its objects before use or recalls a precise earlier setting. Local assumptions, especially mass bounds and positivity of a reference volume, apply only where stated. Proof-local coordinates, maps, measures and parameters are introduced at their first use.
